\begin{prop} \label{prop:limisunique}
A convergent sequence has a unique limit.
\end{prop}

The proof of this proposition exhibits a useful technique in
analysis.  Many proofs follow the same general scheme.  We want to
show a certain quantity is zero.  We write the quantity using the 
triangle inequality as two quantities, and we estimate each one
by arbitrarily small numbers.

\begin{proof}
%NOTE: should be word for word the same as 7.3.3
Suppose $\{ x_n \}_{n=1}^\infty$ has limits $x$ and $y$.
Take an arbitrary $\epsilon > 0$.
From the definition find an $M_1$ such that for all $n \geq M_1$,
$\sabs{x_n-x} < \nicefrac{\epsilon}{2}$.  Similarly, find an $M_2$
such that for all $n \geq M_2$, we have
$\sabs{x_n-y} < \nicefrac{\epsilon}{2}$.
Now take an $n$ such that $n \geq M_1$ and also $n \geq M_2$, and estimate
\begin{equation*}
\begin{split}
\sabs{y-x}
& =
\babs{x_n-x - (x_n -y)} \\
& \leq
\sabs{x_n-x} + \sabs{x_n -y} \\
& <
\frac{\epsilon}{2} + \frac{\epsilon}{2} = \epsilon .
\end{split}
\end{equation*}
As $\sabs{y-x} < \epsilon$ for all $\epsilon > 0$, we have $\sabs{y-x} = 0$
and $y=x$.  Hence, the limit (if it exists) is unique.
\end{proof}
